\documentclass[10pt,a4paper]{amsart}
\usepackage[margin=19mm]{geometry}
\usepackage{amsmath,amssymb,mathtools}
\usepackage[scr=rsfs,cal=euler]{mathalfa}
\usepackage{xfrac}
\usepackage[unicode,hidelinks,bookmarks=false]{hyperref}
\newcommand{\ie}{i.e.\ }
\newcommand{\cf}{cf.\ }
\newcommand{\resp}{respectively\ }
\newcommand{\Subsection}{\subsection}
\providecommand{\nobreakdash}{\nobreak\hbox{-}}
\setlength{\emergencystretch}{2em}
\multlinegap0pt
\usepackage{bm}
\usepackage{bbm}
\usepackage{dsfont}
\usepackage{xspace}
\usepackage{cancel}
\usepackage{mathtools}
\usepackage{amsthm}
\makeatletter
\@ifundefined{theorem}{\newtheorem{theorem}{Theorem}}{}
\@ifundefined{proposition}{\newtheorem{proposition}[theorem]{Proposition}}{}
\makeatother
\newcommand{\norm}[1]{\left\lVert #1 \right\rVert}
\title[Resonant Microstructures as Dirac-type Actuators for Acousti]{Resonant Microstructures as Dirac-type Actuators for Acoustic Wave Control}
\author{Arpan Mukherjee, Mourad Sini}
\date{Source: arXiv:2605.25978v1 (CC BY 4.0)}
\begin{document}
\maketitle

\noindent\emph{Source and licence.} Resonant Microstructures as Dirac-type Actuators for Acoustic Wave Control. Version arXiv:2605.25978v1 (CC BY 4.0), distributed under the Creative Commons Attribution 4.0 International licence, as recorded in the arXiv licence metadata for this exact version and shown on the abstract page https://arxiv.org/abs/2605.25978v1. Full rights notice: licenses/arxiv-2605.25978-CC-BY-4.0.txt.\par\smallskip

\noindent\textit{Editorial scope.} Counted unit: the Proposition whose source label is \texttt{prop:well-posed} in the article (source0.tex lines 1274--1299, source label \texttt{prop:well-posed}), delivered together with the article's own complete original proof of it (source0.tex lines 1301--1337, one \texttt{proof} environment of 37 lines). No mathematics is written, repaired or re-derived: statement and proof are the article's own text line for line. Required context retained uncounted: eq:galerkin-weak (lines 1260--1266); eq:IC-galerkin (lines 1268--1271). Every definition, display and label of those lines is retained unchanged. Omitted from this package and not redistributed: 1--1259, 1267--1267, 1272--1273, 1300--1300, 1338--2720. Nothing inside a retained excerpt is omitted or altered: every retained body is delivered exactly as the article's lines are written, so no omission receipt of any kind accompanies this package. Citations: the retained excerpts carry no citation command at all, so no bibliography entry of the article is transcribed and no reference is redistributed. Reference adaptation: none was needed and none was made; every reference inside the delivered unit resolves inside it, so no unresolved reference or citation is produced. Printed slips kept literal: no printed word, symbol, hypothesis or number of the counted statement or of its proof was altered. Layout and class: the article's own class is \texttt{article} with packages including inputenc, fontenc, flafter, amsmath, amssymb, latexsym, psfrag, graphicx, color, indentfirst, bm, ntheorem, mathrsfs, enumerate; the wrapper is the lane's pinned \texttt{amsart} editorial wrapper at 10pt on A4, which restates the theorem-like environments the retained text needs and adds the article's own macro definitions that the retained text needs (\texttt{\textbackslash norm}), with extra wrapper packages (bm, bbm, dsfont, xspace, cancel, mathtools). The delivered numbering is the unit's own. Assets: no retained span contains a figure environment, a graphics-inclusion command, a PDF-page inclusion, a table or a TikZ picture; no image file is copied into this package, the package declares no asset dependency and no image file is redistributed with it. Licence: version arXiv:2605.25978v1 of this article is distributed under the Creative Commons Attribution 4.0 International licence, as recorded in the arXiv licence metadata for this exact version and shown on the abstract page https://arxiv.org/abs/2605.25978v1. A full rights notice is delivered at licenses/arxiv-2605.25978-CC-BY-4.0.txt. Characters: every retained excerpt is pure ASCII, so the delivered engine, XeLaTeX with its default Latin Modern text font, reports no missing character.
\begin{equation}
\label{eq:galerkin-weak}
 \frac{1}{c_0^2}\langle \partial_t^2 p_M(\cdot,t), \varphi\rangle
 + \langle \nabla p_M(\cdot,t), \nabla \varphi\rangle
 = \sum_{\alpha=1}^N q_\alpha(t)\,\varphi(y_\alpha),
 \qquad \forall \varphi\in H_M,
\end{equation}

\begin{equation}
\label{eq:IC-galerkin}
 p_M(\cdot,0)=0,\qquad \partial_t p_M(\cdot,0)=0.
\end{equation}

\begin{proposition}[Galerkin well-posedness from rest]
\label{prop:well-posed}
Let $T>0$ and $\mathbf q\in L^2(0,T;\mathbb R^N)$. Then the Galerkin problem
\eqref{eq:galerkin-weak}--\eqref{eq:IC-galerkin} admits a unique solution
\[
 p_M\in H^2(0,T;H_M).
\]
Writing $p_M(x,t)=\sum_{k=1}^{N_M} p_k(t)\psi_k(x)$, the coefficients satisfy the decoupled forced oscillator system
\begin{equation}
\label{eq:mode-ODE-N}
 \ddot p_k(t)+\omega_k^2 p_k(t)
 = c_0^2 \sum_{\alpha=1}^N q_\alpha(t)\psi_k(y_\alpha),
 \qquad k=1,\dots,N_M,
\end{equation}
with initial data $p_k(0)=0$ and $\dot p_k(0)=0$.
Moreover, there exists a constant $C_{T,M}>0$ (depending on $T$ and the target subspace $H_M$, but not on $\mathbf{q}$) such that the energy satisfies the bound
\begin{equation}
\label{eq:galerkin-energy-est}
 \sup_{t\in[0,T]}
 \Bigl(
  \norm{\nabla p_M(\cdot,t)}_{L^2(\Omega)}
  + \norm{\partial_t p_M(\cdot,t)}_{L^2(\Omega)}
 \Bigr)
 \le C_{T,M} \norm{\mathbf q}_{L^2(0,T;\mathbb R^N)}.
\end{equation}
\end{proposition}

\begin{proof}
Substituting the expansion $p_M(x,t)=\sum\limits_{k=1}^{N_M} p_k(t)\psi_k(x)$ into the weak formulation \eqref{eq:galerkin-weak} and choosing test functions $\varphi=\psi_k$, we obtain
\[
\frac{1}{c_0^2}\ddot p_k(t) + \lambda_k p_k(t)
= \sum_{\alpha=1}^N q_\alpha(t)\psi_k(y_\alpha).
\]
Multiplying by $c_0^2$ and setting $\omega_k^2=c_0^2\lambda_k$ yields \eqref{eq:mode-ODE-N}. Since $\psi_k\in C^\infty(\overline{\Omega})$, the quantities $\psi_k(y_\alpha)$ are well-defined. 
\newline
Introducing the modal vector $\bm p^M(t):=(p_1(t),\dots,p_{N_M}(t))^\top$, the diagonal matrix $\Lambda_M:=\mathrm{diag}(\omega_1^2,\dots,\omega_{N_M}^2)$, and the geometric modal evaluation matrix $C_M(\mathbf y)\in\mathbb R^{N_M\times N}$ with $[C_M(\mathbf y)]_{k\alpha}:=\psi_k(y_\alpha)$, the system can be written compactly as
\begin{equation}\label{eq:modal-matrix}
\ddot{\bm p}^M(t)+\Lambda_M\bm p^M(t)=c_0^2 C_M(\mathbf y)\mathbf q(t).
\end{equation}
This is a linear second-order ODE in $\mathbb{R}^{N_M}$ with $L^2$ forcing, and therefore admits a unique solution $\bm p^M\in H^2(0,T;\mathbb{R}^{N_M})$ subject to $\bm p^M(0)=0$ and $\dot{\bm p}^M(0)=0$. This implies $p_M\in H^2(0,T;H_M)$. 
\noindent
We now derive the energy estimate. Define the modified modal energy
\[
E_M(t):=\frac12\Bigl(|\Lambda_M^{1/2}\bm p^M(t)|^2+|\dot{\bm p}^M(t)|^2\Bigr).
\]
Differentiating and substituting the dynamics from \eqref{eq:modal-matrix}, we obtain
\[
E_M'(t)=\dot{\bm p}^M(t)\cdot (c_0^2 C_M(\mathbf y)\mathbf q(t)).
\]
By the Cauchy--Schwarz inequality and the matrix operator norm,
\[
E_M'(t)\le c_0^2 |\dot{\bm p}^M(t)|\,\|C_M(\mathbf y)\|\,|\mathbf q(t)|
\le c_0^2 \sqrt{2E_M(t)}\,\|C_M(\mathbf y)\|\,|\mathbf q(t)|.
\]
Setting $y(t):=\sqrt{E_M(t)}$ and noting $y(0)=0$, we obtain $y'(t)\le c_0^2 \|C_M(\mathbf y)\|\,|\mathbf q(t)|$ for a.e.\ $t\in(0,T)$. Integrating over $[0,t]$ yields
\[
\sup_{t\in[0,T]} y(t) \le c_0^2 \|C_M(\mathbf y)\|\sqrt{T}\,\|\mathbf q\|_{L^2(0,T;\mathbb R^N)}.
\]
Finally, noting that $|\dot{\bm p}^M(t)| = \|\partial_t p_M(\cdot,t)\|_{L^2(\Omega)}$ and translating the stiffness term via $\omega_k^2 = c_0^2\lambda_k$:
\[
|\Lambda_M^{1/2}\bm p^M(t)|^2 = \sum_{k=1}^{N_M} c_0^2\lambda_k |p_k(t)|^2 = c_0^2 \|\nabla p_M(\cdot,t)\|_{L^2(\Omega)}^2,
\]
the bounds on $y(t)$ directly establish \eqref{eq:galerkin-energy-est}. \hfill $\blacksquare$
\end{proof}

\end{document}
